% HYPERGEOMCENTER1
\documentclass[11pt]{article}
\usepackage[margin=.88in]{geometry}
\usepackage{amsmath,amssymb,amsthm}
\newtheorem{proposition}{Proposition}
\begin{document}
\begin{center}
{\Large HYPERGEOMCENTER1: the centered cross moment has no product ratio}\\[4pt]
\textbf{Date: 2026-10-04 22:15 PT}\\
\textbf{Status: PROVED/EXACT}\\
\textbf{P versus NP status: OPEN}
\end{center}

\section{What was already committed}
SLICEWIDTH1 is on \texttt{main} as \texttt{912d97a0}. This note continues from HYPERGEOM-MOMENT1, whose conjunction ratio is proved and whose centered-moment step was left as a candidate.

\section{The ratio checks}
DISJOINT-SLICE-RATIO says $1-R_{a,b}\le ab/(t-a-b+1)$. Direct evaluation at density $3.5$ holds on every tested block:
\[
\begin{array}{c|c|c|c}
q&(a,b)&1-R&\text{bound}\\
\hline
20&(1,1)&0.01418&0.01449\\
20&(2,2)&0.05631&0.05970\\
20&(3,3)&0.12400&0.13846\\
20&(6,6)&0.42592&0.61017
\end{array}
\]
The same inequality holds at $q=16$ and $q=24$.

\section{Centered moments do not factor}
\begin{proposition}[CENTER-RATIO-FAIL]
Let $A,B$ be nonempty disjoint coordinate blocks with $a+b\le t$. Under the uniform exact-$t$ slice,
\[
\mathbb E[1_A 1_B]-\mathbb E[1_A]\mathbb E[1_B]
=\mathbb E[1_A]\mathbb E[1_B](R_{a,b}-1)<0.
\]
Under any product measure the same centered cross moment is identically $0$. Therefore no multiplicative distortion bound of the form
\[
\text{slice centered moment}=(1+o(1))\times\text{product centered moment}
\]
can hold: the denominator is zero and the numerator is not.
\end{proposition}
\begin{proof}
The identity is DISJOINT-SLICE-RATIO. Each factor of $R_{a,b}$ is strictly less than one for $1\le a\le t-b$, so $R_{a,b}<1$. A product measure has independent coordinates, so its cross cumulant of order $(a,b)$ vanishes.
\end{proof}

The additive gap is small because both probabilities are small. At $q=20$ it is $8.35\times 10^{-7}$ for $(1,1)$ and $2.33\times 10^{-14}$ for $(3,3)$. Small additive error is not a ratio bound.

\section{Status}
CENTER-RATIO-FAIL: \textbf{PROVED/EXACT}. It kills the ratio form of the HYPERGEOM-PARTITION candidate. The additive window $\exp(O((pw)^2/t))$ remains unproved, and centering can still amplify it. The lemma is generic to every fixed-weight slice, including the GE-easy XOR slice, so it is not a SAT obstruction. $P\ne NP$ remains OPEN.
\end{document}
